\section{Information-State Transfer, Calibration, and Control}
\label{sec:transfer_operational_implications}

The horizon law is attached to an information state, not only to a raw
retention window. Once a pipeline admits a finite-sample envelope and a
staleness envelope, the same constitutive balance governs estimation,
calibration, and control. Finite projection families and sliced geometries
transfer the law under uniform projected control. Beyond those proved
classes, high-dimensional practicality becomes a geometric question: temporal
validity remains inferable only when the chosen state geometry preserves a
usable finite-sample exponent.

\begin{corollary}[Effective-exponent horizon transfer]
\label{cor:effective_exponent_horizon_transfer}
Let $d_{\mathrm{eff}}$ be a probability metric such that, on the drifting class of
interest,
\[
  \mathbb E\,d_{\mathrm{eff}}\!\left(\widehat P_t^{(n)},\bar P_t^{(n)}\right)
  \le C_{K,\mathrm{eff}} n^{-a_{\mathrm{eff}}}
\]
for some $a_{\mathrm{eff}}>0$, and suppose the same class satisfies the staleness
bound
\[
  d_{\mathrm{eff}}\!\left(\bar P_t^{(n)},P_t\right)
  \le C_{S,\mathrm{eff}}\zeta n^H.
\]
Then
\[
  \mathbb E\,d_{\mathrm{eff}}\!\left(\widehat P_t^{(n)},P_t\right)
  \le C_{K,\mathrm{eff}} n^{-a_{\mathrm{eff}}} + C_{S,\mathrm{eff}}\zeta n^H,
\]
and the induced validity horizon satisfies
\[
  n^*_{\mathrm{eff}}(H)
  \asymp
  (C_{K,\mathrm{eff}}/\zeta)^{1/(a_{\mathrm{eff}}+H)}.
\]
\end{corollary}

\begin{proof}
Apply \Cref{thm:abstract_upper_law} with $d=d_{\mathrm{eff}}$.
\end{proof}

The relevant question is which geometry preserves a usable exponent
$a_{\mathrm{eff}}$. Finite projection and sliced constructions provide explicit
examples with one-dimensional finite-sample control. Intrinsic-dimension and
projection-robust transport can replace ambient dimension by an effective
dimension, while smoothed transport can preserve parametric behavior through
regularization \citep{NilesWeedBach2019,PatyCuturi2019,NietertEtAl2022}. In
these broader settings, the effective-exponent corollary yields a horizon law;
it does not imply a new multivariate transport theorem.

A practical guide is therefore geometry-first. When the ambient dimension is low
and the geometry is well resolved, raw Wasserstein remains the most direct
choice. When only a small set of directions carries drift, finite projections or
sliced variants are preferable because they preserve a usable exponent by
discarding uninformative directions. When the geometry is noisy or the support
is moderate, a calibrated regularization band is natural because it regularizes
the finite-sample term without forcing a new transport theorem. In all cases,
the guiding criterion is whether the chosen geometry supplies an effective
exponent large enough to keep the horizon nondegenerate.

An equivalent operator view is often more informative than the geometry-first
view. Let $R_t^{(n)}$ retain a state from the recent past, let $U_t$ update that
state as new observations arrive, and let $G_t$ extract an estimate, interval,
alarm statistic, or control action. The object governed by validity is
the composite information state $G_t\circ U_t\circ R_t^{(n)}$, not memory length
alone. The horizon law explains when a retained state remains useful, the
inference layer estimates that horizon from lag geometry, and the control layer
decides whether to deploy a point estimate, a useful-memory band, or a deferred
update.

Conformal calibration under drift fits this template directly. A split-conformal
calibration set is a retained information state, the nonconformity quantile is a
readout, and predictive coverage is the operational validity criterion. The same
finite-sample versus staleness balance therefore applies to calibration-window
selection. In the benchmark reported here, the calibration score has an interior
optimum over window size, coverage degrades before alarm, and the empirically
safe calibration band coincides with the useful-memory band. The conformal
connection is therefore constitutive rather than analogical: calibration under
drift is another readout of validity.

Sequential memory control closes the same object on the action side. Once memory
becomes a controlled state and updates incur fixed and proportional actuation
cost, the decision problem is not to maximize memory or to react as early as
possible, but to remain inside the temporally valid region at acceptable cost.
The resulting frontier is regime dependent. Drift roughness, sensing quality,
and actuation price determine whether aggressive control, detector-triggered
updates, or more conservative deployment is preferred.

Weighted-window diagnostics on uniform, triangular, and geometric tapers show
the same finite-sample versus lag trade-off in operational form; see
\Cref{tab:weight_sensitivity_summary}. Horizon inference begins with the plug-in
horizon entering the useful-memory region and remaining there under streaming
noise. Support change beyond the fixed support-growth construction and online
validity control under streaming noise are not treated here at theorem level.
